\documentclass[11pt,a4paper]{article}
\usepackage[margin=2.1cm]{geometry}
\usepackage[T1]{fontenc}
\usepackage[utf8]{inputenc}
\usepackage{booktabs}
\usepackage{array}
\usepackage{tabularx}
\usepackage{xcolor}
\usepackage{hyperref}
\usepackage{listings}
\usepackage{float}

\hypersetup{colorlinks=true, linkcolor=black, urlcolor=blue}
\setlength{\parindent}{0pt}
\setlength{\parskip}{0.55em}
\emergencystretch=2em

\lstset{
basicstyle=\ttfamily\small, breaklines=true, frame=single, columns=fullflexible,
keepspaces=true, showstringspaces=false }

\title{STM32N6 Square Segmentation Workflow\\Host Preparation, Cube.AI Project Generation, and UART Run Guide}
\author{Project documentation}
\date{June 23, 2026}

\begin{document}
\maketitle

\tableofcontents
\newpage

\begin{abstract}
This document describes the STM32N6 workflow used for binary
square segmentation on grayscale 128x128 images. The workflow covers host-side ONNX
export, INT8 quantization, STM32Cube.AI project generation, firmware customization for
UART-based inference, signed flashing to external XSPI, and host-to-board inference
through a binary protocol.

The repository is split into two main parts. The host-side model preparation flow lives
under \path{models/square_seg_stm32n6/}. The board-side firmware project generated from
the INT8 model and then customized for the final deployment flow lives under
\path{stm32-model-example-int8/}.
\end{abstract}

\section{Project Overview}

The goal of this project is binary square segmentation on STM32N6 with a UART-based
host/board inference loop. A host script sends one quantized 128x128 grayscale image to
the board, the STM32N6 runs inference with the generated AI runtime, and the board
returns the predicted mask together with the measured inference time.

The key repository paths are:

\begin{table}[H]
\centering
\caption{Main repository areas.}
\begin{tabular}{p{0.32\linewidth} p{0.58\linewidth}}
\toprule
Path & Role \\
\midrule
\path{models/square_seg_stm32n6/configs/} & YAML presets for export, quantization, compare, deployment, and UART inference \\
\path{models/square_seg_stm32n6/src/} & Host-side Python scripts \\
\path{models/square_seg_stm32n6/artifacts/} & Generated ONNX, INT8, report, and UART outputs \\
\path{stm32-model-example-int8/} & STM32Cube.AI generated project, later customized for the final UART workflow \\
\path{stm32-model-example-int8/tools/deploy_signed.sh} & Signed flashing helper for FSBL and Appli \\
\bottomrule
\end{tabular}
\end{table}

\section{Model Preparation Pipeline}

The host-side flow that leads to the deployable model is kept in
\path{models/square_seg_stm32n6/}. The default configuration in
\path{configs/train.yaml}, uses grayscale 128x128 images, and defines a compact
square segmentation CNN with 32 channels and 13 layers.

\subsection{Python environment}

From the repository root:

\begin{lstlisting}[language=bash]
python3 -m venv .venv
source .venv/bin/activate
pip install -r requirements.txt
\end{lstlisting}

\subsection{Export ONNX}

The ONNX export step is implemented by
\path{models/square_seg_stm32n6/src/export_onnx.py} using the hardcoded configuration file
\path{models/square_seg_stm32n6/configs/export.yaml}.

\begin{lstlisting}[language=bash]
python models/square_seg_stm32n6/src/export_onnx.py
\end{lstlisting}

This step loads \path{artifacts/model.pt}, rebuilds the model from the saved training
configuration, exports a fixed-shape ONNX graph to
\path{models/square_seg_stm32n6/artifacts/model.onnx}, validates the ONNX graph with
\texttt{onnx.checker}. By default it also runs an ONNX Runtime CPU check and writes
\path{models/square_seg_stm32n6/artifacts/onnx_report.json}.

\subsection{Quantize to INT8}

The quantization step is implemented by
\path{models/square_seg_stm32n6/src/quantize_onnx.py} using the hardcoded configuration file
\path{models/square_seg_stm32n6/configs/quantize.yaml}.

\begin{lstlisting}[language=bash]
python models/square_seg_stm32n6/src/quantize_onnx.py
\end{lstlisting}

This step generates \path{models/square_seg_stm32n6/artifacts/model_int8.onnx} and
writes \path{models/square_seg_stm32n6/artifacts/quantization_report.json}. The report
shows the deployment constraints that remain important later:

\begin{itemize}
\item the source ONNX must be exported with opset 13 or newer; this repository uses opset
      13,
\item static calibration uses 256 deterministic samples with seed 42,
\item calibration preprocessing is grayscale, 128x128, and divide-by-255,
\item the quantized graph stays in QDQ INT8 format.
\end{itemize}

\section{STM32Cube.AI Generation}

The committed STM32 firmware project was generated in STM32Cube.AI from the INT8 model,
then customized manually afterward. In practice, the INT8 model imported into the STM32
tooling is the same final quantized model produced by the host-side flow, namely
\path{models/square_seg_stm32n6/artifacts/model_int8.onnx}.

STM32Cube.AI generated the project \path{stm32-model-example-int8/}. The resulting
project includes the standard STM32N6 secure boot split between an FSBL context and an
application context, together with the generated AI runtime artifacts used by the final
firmware.

The most important generated directories are:

\begin{table}[H]
\centering
\caption{Cube.AI generated project structure.}
\begin{tabular}{p{0.28\linewidth} p{0.62\linewidth}}
\toprule
Path & Meaning \\
\midrule
\path{stm32-model-example-int8/.ai/} & Cube.AI internal workspace, reports, generated network package, and auxiliary output files \\
\path{stm32-model-example-int8/Appli/AI/App/} & Application-integrated AI runtime sources such as \texttt{network.c}, \texttt{stai\_network.c}, and \texttt{network\_atonbuf.xSPI2.c} \\
\path{stm32-model-example-int8/FSBL/} & First-stage bootloader sources for the secure external-memory boot flow \\
\path{stm32-model-example-int8/STM32CubeIDE/} & IDE workspace containing separate \texttt{FSBL} and \texttt{Appli} projects and their build outputs \\
\bottomrule
\end{tabular}
\end{table}

The \path{.ai/} directory shows the typical generated content that came from
STM32Cube.AI, including:

\begin{itemize}
\item \path{.ai/st_ai_output/network.c},
\item \path{.ai/st_ai_output/network.h},
\item \path{.ai/st_ai_output/network_atonbuf.xSPI2.c},
\item \path{.ai/st_ai_output/stai_network.c},
\item \path{.ai/st_ai_output/network_generate_report.txt},
\item \path{.ai/st_ai_output/network_model_int8_c_info_generate.json}.
\end{itemize}

The project metadata in \path{stm32-model-example-int8/stm32-model-example-int8.ioc}
shows the secure multi-context structure with \texttt{Mcu.Context0=FSBL},
\texttt{Mcu.Context1=Appli}, and \texttt{Mcu.Context2=ExtMemLoader}. After generation,
the project was manually adapted to support the final end-to-end UART inference flow.

\section{Firmware Integration After Cube.AI Generation}

The final custom behavior is implemented in
\path{stm32-model-example-int8/Appli/Core/Src/main.c}. The documentation here focuses on
the final behavior rather than the edit history.

At startup, the application:

\begin{itemize}
\item initializes the STM32N6 AI runtime,
\item resolves the generated input and output tensor buffers,
\item enables DWT cycle counting for board-side timing,
\item waits forever for frames on \texttt{USART1}.
\end{itemize}

The application then exposes a simple inference service over UART:

\begin{itemize}
\item a host sends a \texttt{HELLO} request,
\item the board returns image and quantization metadata,
\item the host sends one CRC-protected quantized input tensor,
\item the board runs inference on the generated network,
\item the board returns the execution time in microseconds and the raw INT8 output tensor,
\item error conditions are mapped to protocol status codes and returned as \texttt{ERROR}
      frames.
\end{itemize}

The same application source also contains the board-side CRC32 logic, frame parsing,
quantized input reception, timing capture, and response serialization. The output mask
returned to the host is the raw INT8 network output. Thresholding and metric computation
are intentionally done on the host side so the run script can save both the prediction
artifact and the evaluation report.

The project layout also reflects the final XSPI boot and deployment arrangement. The
board image is split into trusted FSBL and trusted application binaries, and the
deployment helper flashes them separately into external XSPI using the dedicated
STM32N6570-DK external loader.

\section{UART Protocol and Host Workflow}

\subsection{Protocol shape}

The host script \path{models/square_seg_stm32n6/src/uart_inference.py} and the firmware
in \path{stm32-model-example-int8/Appli/Core/Src/main.c} implement the same protocol.

Every frame starts with a 16-byte header:

\begin{table}[H]
\centering
\caption{UART frame header.}
\begin{tabular}{ll}
\toprule
Field & Meaning \\
\midrule
4 bytes & magic \texttt{SSEG} \\
1 byte & protocol version, fixed to \texttt{1} \\
1 byte & message type \\
2 bytes & status code, little-endian \\
4 bytes & payload length, little-endian \\
4 bytes & CRC32 of payload, or 0 for empty payload \\
\bottomrule
\end{tabular}
\end{table}

The message types are:

\begin{itemize}
\item \texttt{HELLO} (\texttt{0x01}),
\item \texttt{IMAGE} (\texttt{0x02}),
\item \texttt{HELLO\_ACK} (\texttt{0x81}),
\item \texttt{RESULT} (\texttt{0x82}),
\item \texttt{ERROR} (\texttt{0xFF}).
\end{itemize}

The \texttt{HELLO\_ACK} payload is a 24-byte metadata block with:

\begin{itemize}
\item input width and height,
\item input tensor size in bytes,
\item output tensor size in bytes,
\item input scale and input zero-point,
\item output scale and output zero-point.
\end{itemize}

The \texttt{RESULT} payload contains:

\begin{itemize}
\item 8 bytes of inference time in microseconds,
\item the raw INT8 output tensor.
\end{itemize}

The firmware uses explicit error status codes for bad magic, bad protocol version, bad
message, bad payload length, bad CRC, AI initialization failure, AI inference failure,
and UART transfer failure.

\subsection{Host workflow}

The default host configuration is
\path{models/square_seg_stm32n6/configs/uart_inference.yaml}. It defines:

\begin{itemize}
\item serial port \texttt{/dev/ttyACM0},
\item baud rate \texttt{115200},
\item sample image \path{../../../../data/data/test/img_1981_image.tif},
\item reference mask \path{../../../../data/data/test/img_1981_label.tif},
\item output directory \path{models/square_seg_stm32n6/artifacts/uart_inference/}.
\end{itemize}

When the user runs:

\begin{lstlisting}[language=bash]
python models/square_seg_stm32n6/src/uart_inference.py
\end{lstlisting}

the host script performs the following final workflow:

\begin{enumerate}
\item open the serial port and keep retrying until the board answers the handshake,
\item send \texttt{HELLO},
\item receive tensor metadata in \texttt{HELLO\_ACK},
\item load the configured grayscale image,
\item resize it if needed and quantize it with the board-provided input scale and
      zero-point,
\item send the \texttt{IMAGE} payload with CRC32 protection,
\item receive \texttt{RESULT},
\item decode the returned inference time and raw INT8 mask,
\item threshold the mask with the returned output zero-point,
\item compare the predicted binary mask with the reference mask,
\item save a PNG mask and a JSON report.
\end{enumerate}

The generated files are:

\begin{itemize}
\item \path{models/square_seg_stm32n6/artifacts/uart_inference/prediction.png},
\item \path{models/square_seg_stm32n6/artifacts/uart_inference/report.json}.
\end{itemize}

The JSON report includes:

\begin{itemize}
\item inference time in microseconds and milliseconds,
\item Dice score,
\item IoU score,
\item protocol magic and version,
\item tensor dimensions,
\item tensor sizes,
\item input and output quantization parameters.
\end{itemize}

\section{Build and Deployment}

\subsection{Build the firmware}

Open the CubeIDE workspace under \path{stm32-model-example-int8/STM32CubeIDE/}. Build
both projects:

\begin{itemize}
\item \texttt{FSBL},
\item \texttt{Appli}.
\end{itemize}

The relevant output binaries are the latest:

\begin{itemize}
\item \path{stm32-model-example-int8/STM32CubeIDE/FSBL/Debug/stm32-model-example-int8_FSBL.bin},
\item \path{stm32-model-example-int8/STM32CubeIDE/Appli/Debug/stm32-model-example-int8_Appli.bin}.
\end{itemize}

The deploy helper locates the newest matching \texttt{FSBL.bin} and \texttt{Appli.bin}
automatically, so the normal user command does not need explicit file paths.

\subsection{Signed flashing flow}

The final deployment command is:

\begin{lstlisting}[language=bash]
bash stm32-model-example-int8/tools/deploy_signed.sh
\end{lstlisting}

This script performs the final clean flow:

\begin{enumerate}
\item locate the latest FSBL and Appli \texttt{.bin} files under \path{STM32CubeIDE/},
\item add trusted STM32N6 headers with \texttt{STM32\_SigningTool\_CLI},
\item use \path{MX66UW1G45G_STM32N6570-DK.stldr} as the external loader,
\item flash the trusted FSBL to \texttt{0x70000000},
\item flash the trusted Appli to \texttt{0x70100000},
\item reset the target.
\end{enumerate}

The script resolves \texttt{STM32\_Programmer\_CLI}, \texttt{STM32\_SigningTool\_CLI},
and the external loader automatically when possible. The default programmer connect
arguments are:

\begin{lstlisting}
port=SWD mode=HOTPLUG
\end{lstlisting}

The required board state during programming is:

\begin{itemize}
\item DEV boot mode,
\item BOOT1 high.
\end{itemize}

If \path{stm32-model-example-int8/out/deploy} is not writable, the script falls back
automatically to:

\begin{lstlisting}
/tmp/stm32n6-deploy-$USER
\end{lstlisting}

and stores the trusted binaries there.

\section{Run Guide}

The clean final user workflow is:

\begin{enumerate}
\item prepare the Python environment,
\item export the ONNX model,
\item quantize it to INT8,
\item open \path{stm32-model-example-int8/STM32CubeIDE/} in STM32CubeIDE,
\item build \texttt{FSBL} and \texttt{Appli},
\item run \texttt{bash stm32-model-example-int8/tools/deploy\_signed.sh},
\item reset or power-cycle the board if needed,
\item run \texttt{python models/square\_seg\_stm32n6/src/uart\_inference.py}.
\end{enumerate}

In command form:

\begin{lstlisting}[language=bash]
python3 -m venv .venv
source .venv/bin/activate
pip install -r requirements.txt

python models/square_seg_stm32n6/src/export_onnx.py

python models/square_seg_stm32n6/src/quantize_onnx.py

bash stm32-model-example-int8/tools/deploy_signed.sh

python models/square_seg_stm32n6/src/uart_inference.py
\end{lstlisting}

The produced outputs for one run are:

\begin{itemize}
\item a predicted mask PNG,
\item a JSON report with inference time, Dice, IoU, and tensor metadata.
\end{itemize}

That is the final workflow represented by this repository: host-side preparation of the
INT8 model, STM32Cube.AI generation of the STM32 project from that INT8 model, manual
firmware customization for the UART service, signed flashing to XSPI, and host-to-board
inference through the defined serial protocol.

\end{document}
